\section{Finite graphs with exactly uniform demands}\label{sec:finite-graph}

We now convert Theorem~\ref{thm:score-gap} into the graph promised in
Theorem~\ref{thm:main}.  The conversion must retain cuts of arbitrarily small
positive mass.  In particular, replacing the demand by an additive approximation
alone would be insufficient.  We keep a multiplicative upper bound on the
retained demand and attach every remaining vertex to that demand.

Fix an integer $C'>C$, and set
\[
 A_0=10^{12},\qquad \kappa=C'A_0,\qquad T_0=100\kappa.
\]
Apply Theorem~\ref{thm:score-gap} with this $T_0$.  Write $\mu$ for its demand
law, and $\cD$ for the sum of its raw comparison forms.  Thus, for a perfect
witness $x$, the total truth cost is at most $4$, and
\begin{equation}\label{eq:demand-tails-finite}
 \Prob\{B^0(x)\ge1/4\}\ge10^{-4},\qquad
 \Prob\{B^0(x)\le-1/4\}\ge10^{-4}.
\end{equation}
In a NO instance, every measurable binary lookup $h$ with
$\nu=\Var_\mu(h)>0$ satisfies $\cD(h)>T_0\nu$.
All arities, numerical parameters, and total raw comparison weights below
are constants depending on $C$ only.

\subsection{Canonical keys and rational integration}

Let $b$ be the number of bit names in the test system. After fixing all
parameters of the score construction, unroll each complete demand or
comparison experiment, including both endpoints. Count all primitive
query occurrences: retained, frozen, and future primary lists, resampled
copies, fresh and forced path lists, and all lists defining formal
partitions. Counting repetitions gives an effectively computable bound
$Q_*=Q_*(C)$ on the number of distinct queried bit names in any experiment.
Descriptors, cell indicators, translations, and numerical noises introduce
no names beyond those lists. The finite sampler templates and their
sample counts are fixed by $C$, so $Q_*$ is independent of input size.
Choose $k=Q_*+1$; here $k$ is solely the key-enumeration bound.
Choose a bounded rational interval containing the ranges of all score
functions in these experiments, and a fixed finite rational grid in that
interval. Rounding means nearest-grid-point rounding
with a fixed rule at ties.  A \emph{key} is a grid-valued function on the bit
cube, specified by its ordered list of essential bit names and its truth
table.  Equal functions have the same key; sampling presentations are not
part of a key.  Enumerate all such keys of arity at most $k$.  If the grid
has $g$ values, there are at most
\[
 \sum_{j=0}^k\binom bj g^{2^j}
\]
keys.  This is polynomial in the original instance size.

Write $R(B)$ for entrywise rounding of a score.  If the entrywise error is at
most $\delta$ and the total raw comparison weight is $W_0$, then for each
truth assignment
\[
 \abs{R(B)(x)-R(\widetilde B)(x)}
 \le \abs{B(x)-\widetilde B(x)}+2\delta.
\]
Take the fixed grid fine enough that $2\delta W_0\le1$ and $\delta<1/16$.
The total rounded truth cost is then at most $5$, and the two demand tails
remain strictly outside $[-1/8,1/8]$.

Every coloring of the keys extends to a measurable lookup on the scores:
apply $R$ and then the coloring.  The range can be truncated outside the
interval used by the experiments and the function defined arbitrarily on
any unused scores.  Measurability follows because each rounded table is
specified by finitely many polynomial inequalities in the numerical
parameters of each experiment.  Consequently the soundness assertion of
Theorem~\ref{thm:score-gap} applies to every key coloring.

Here and below, $\mu_i$ denotes the pushforward demand mass of key $i$.
For an unordered pair of distinct keys, let $c^0_{ij}$ be its total raw
comparison weight, adding both possible orderings.  Self-comparisons
contribute zero.  Thus
\begin{equation}\label{eq:raw-capacities}
 \cD(h)=\sum_{i<j}c^0_{ij}\abs{h_i-h_j}.
\end{equation}
The numbers $c^0_{ij}$ and $\mu_i$ need not be rational.  The next observation
gives the particular approximations we require without exact real integration.

\begin{lemma}[Rational compilation]\label{lem:rational-compilation}
In polynomial time one can compute nonnegative rational numbers $c_{ij}$
and $l_i$ such that
\begin{align}
 &c_{ij}\ge c^0_{ij},\qquad
   \sum_{i<j}(c_{ij}-c^0_{ij})\le1,\label{eq:capacity-upper}\\
 &0\le l_i\le\mu_i,\qquad L:=\sum_i l_i\ge1-10^{-6}.
   \label{eq:demand-lower}
\end{align}
Their binary encoding lengths are polynomial.  Every positive $l_i$ is at
least an inverse polynomial in the input size, with the polynomial depending
only on $C$.
\end{lemma}

\begin{proof}
Enumerate the discrete seeds in the demand and comparison experiments.
There are polynomially many seeds: the number of samples per experiment is
fixed, and every bit-test sampling space has polynomial size.  Their
probabilities are rational with polynomial encoding length, and every
positive discrete sampling probability is at least an inverse polynomial.

Conditional on a seed, the remaining numerical variables range over
rectangular boxes of fixed dimension.  Atoms and signs are separate discrete
choices.  The smoothing densities and comparison weights are polynomial
with rational coefficients on finitely many rectangular pieces; absolute
derivative weights are split at their changes of sign.  Each score-table
entry is a polynomial with fixed rational coefficients in these variables.
These assertions also hold for forced-slot comparisons, common-noise
comparisons, path interpolation, cell translation, and the formal cell sums:
all involve a fixed number of additions and products of the prescribed
amplitudes, marks, and shifts.  Only the local Boolean tables of the queried
predicates and their bit-name identifications enter these formulas.  Their
sizes are bounded, so, after forgetting the actual bit names, only a fixed
finite collection of numerical patterns occurs.  In particular, the
lookup-dependent reference sign used in the proof of soundness is not used
to define any comparison weight. The input-dependent rational seed
probabilities remain outside these conditional numerical integrals;
they do not enlarge the finite family of polynomial coefficients.

Subdivide each numerical box into rational boxes.  For every grid boundary
and every table entry, first check whether the polynomial entry is
identically equal to that boundary; if so, resolve the tie by the fixed
rounding rule.  Otherwise its equality set has Lebesgue measure zero.
Indeed, a nonzero polynomial in one variable has finitely many roots; in
higher dimension expand in the last variable, apply the induction
hypothesis to the simultaneous vanishing of its coefficients, and then
apply the one-variable assertion on the remaining fibers.  Integration
over fibers proves the claim.

Coefficient bounds on a bounded box give a rational Lipschitz bound for
every polynomial entry.  Comparing its value at a rational box corner
with this bound times the mesh identifies every box on which rounding
is determined.  Apart from the finitely many null boundary sets, each
point is eventually in such a box.  Since the weights are bounded, the
total weighted mass of unresolved boxes tends to zero.

On a resolved box integrate its polynomial weight term by term, obtaining
an exact rational mass.  For demand, discard unresolved boxes.  For
capacities, assign the full box mass to every table pair that could occur
there. One may conservatively use every grid-valued table on the queried
coordinates: each has at most $Q_*<k$ essential names and hence a key in
the enumerated universe. Their number is a constant. This gives an upper bound for
each pair mass.  Its excess is bounded by that constant times the weighted
unresolved mass.  Take the mesh fine enough that demand loses at most
$10^{-6}$ and capacity gains at most $1$ in total.  Because the numerical
patterns form a fixed finite collection, one fixed mesh works for every
seed and every input.  Equivalently, rational subdivision can continue
until the explicit unresolved-mass bounds meet these tolerances.  Averaging
these bounds with the seed probabilities gives the stated errors; the
errors are not summed without those probabilities. More explicitly, let
$U_{f,\omega,\ell}$ be the unresolved polynomial mass of rectangular
piece $\ell$ in comparison form $f$, conditional on seed $\omega$,
with its external form weight included. If $\pi_{f,\omega}$ is the seed
probability and $R_{f,\omega,\ell}$ bounds the number of candidate
unordered distinct key pairs, require
\[
 \sum_f\sum_\omega\pi_{f,\omega}
       \sum_\ell R_{f,\omega,\ell}U_{f,\omega,\ell}\le1.
\]
All formal-cell, derivative-coordinate, and rectangular pieces are
counted with their multiplicities. The analogous demand sum, with no
pair factor, must be at most $10^{-6}$. These are exact rational stopping
tests. A common mesh depth succeeds for every numerical pattern; an
input that stops earlier uses a depth no larger than that fixed bound.

At this fixed mesh, the positive retained conditional demand-box masses
belong to a fixed finite set of positive rational numbers. The same is
true over all mesh depths up to the fixed bound. Each therefore
has a positive constant lower bound.  A positive $l_i$ includes one of
these masses multiplied by a positive discrete seed probability, proving
the inverse-polynomial lower bound.  There are polynomially many seeds
and a constant number of boxes and table entries per seed.  Their rational
arithmetic, including the final sums, has polynomial bit complexity.
\end{proof}

Normalize the retained demand by $\widehat\mu_i=l_i/L$.  From
Equations~\eqref{eq:demand-lower} we have, for every key and every set $E$ of
keys,
\begin{equation}\label{eq:retained-demand}
 \widehat\mu_i\le2\mu_i,\qquad
 \widehat\mu(E)\ge\mu(E)-10^{-6}.
\end{equation}
For the second assertion use $\widehat\mu(E)\ge l(E)$ and
$\mu(E)-l(E)\le1-L$.

\subsection{Replication and auxiliary vertices}

Let $m$ be the number of keys, and choose an integer $K$ at least
$4\kappa m$ and at least the reciprocal of the smallest positive
$\widehat\mu_i$.  Lemma~\ref{lem:rational-compilation} permits $K$ of
polynomial magnitude.  Replace key $i$ by a cluster of
\[
 n_i=\begin{cases}
 \lceil K\widehat\mu_i\rceil,&\widehat\mu_i>0,\\
 1,&\widehat\mu_i=0
 \end{cases}
\]
vertices, designating one representative.  Put $N=\sum_i n_i$ and
$w_i=n_i/N$.  Since $K\le N\le K+m$, we obtain
\begin{equation}\label{eq:replication-weights}
 \frac12\widehat\mu_i\le w_i\le2\widehat\mu_i
 \quad(\widehat\mu_i>0),\qquad
 \alpha:=\sum_{\widehat\mu_i=0}w_i\le\frac1{4\kappa}.
\end{equation}
For the upper bound use $K\widehat\mu_i\ge1$; for the lower bound use
$K\ge m$.  Every key has a positive number of copies.

Place capacity $c_{ij}$ between representatives $i,j$.  Join every other
vertex in cluster $i$ to its representative by an edge of capacity
$2\kappa$.  Finally, if $\widehat\mu_i=0$ and $\widehat\mu_j>0$, add capacity
\[
 2\kappa w_i\widehat\mu_j
\]
between their representatives.  These last edges are the \emph{tethers}.
Add capacities when edges coincide, and set all diagonal capacities to
zero.  Give every unordered pair of distinct vertices demand exactly one.
The output threshold is the positive rational number
\begin{equation}\label{eq:final-threshold}
 a=\frac{A_0}{N^2}.
\end{equation}

\subsection{Completeness}

Suppose $x$ is a perfect witness.  For a common threshold
$\theta\in[-1/8,1/8]$, color key $i$ by
$h_i=\ind_{\{i(x)>\theta\}}$, and give its copies the same color.
The tails in Equation~\eqref{eq:demand-tails-finite}, their preservation by
rounding, and Equations~\eqref{eq:retained-demand}--\eqref{eq:replication-weights}
show that both sides of every such cut have counting mass at least
\[
 \tfrac12(10^{-4}-10^{-6})\ge\tau,\qquad
 \tau:=10^{-4}/4.
\]
Thus all these cuts are proper and nonempty.

Take $\theta$ uniformly in that interval.  For any two real values $u,v$,
the probability that the threshold separates them is at most $4\abs{u-v}$.
The expected raw rounded comparison capacity is consequently at most
$4\cdot5=20$.  The excess in Equation~\eqref{eq:capacity-upper} contributes
at most $1$, and the total tether capacity is
$2\kappa\alpha\le1/2$.  No coherence edge crosses.  In particular, the
expected cut capacity is at most $22$, so some threshold cut has capacity
at most $22$.  If its counting mass is $z$, then $z,1-z\ge\tau$ and
\[
 \frac{\operatorname{cap}(S,V\setminus S)}{|S|(N-|S|)}
 \le\frac{22}{N^2\tau^2}
 <\frac{A_0}{N^2}=a,
\]
because $A_0\tau^2=625>22$.

\subsection{Soundness}

Suppose the formula is unsatisfiable and a nonempty proper cut has ratio
at most $C'a$.  Let $z=|S|/N$ and let $T$ be its capacity.  Then
\begin{equation}\label{eq:putative-small-cut}
 T\le\kappa z(1-z)\le\kappa/4.
\end{equation}
A split cluster would cut an edge of capacity $2\kappa$, so every cluster
is monochromatic.  The cut therefore determines a binary coloring $h_i$
of the keys and hence a measurable lookup on scores.  Set
\[
 \nu=\Var_\mu(h),\qquad
 L_{\rm aux}=\sum_{\widehat\mu_i=0}w_i
                    \sum_j\widehat\mu_j\abs{h_i-h_j}.
\]
The tethers contribute exactly $2\kappa L_{\rm aux}$.

We claim
\begin{equation}\label{eq:variance-transfer}
 z(1-z)\le36\nu+L_{\rm aux}.
\end{equation}
To prove this, draw a key with distribution $w$.  Retain it if its
$\widehat\mu$ mass is positive, and otherwise replace it by an independent
$\widehat\mu$ draw.  The color changes with probability $L_{\rm aux}$.
The replacement distribution is
\[
 \rho_j=\ind_{\{\widehat\mu_j>0\}}w_j+\alpha\widehat\mu_j
 \le3\widehat\mu_j\le6\mu_j.
\]
Apply this coupling independently to two $w$ draws.  Their original
disagreement probability is $2z(1-z)$; their replacement disagreement
probability is at most $36\cdot2\nu$ by the product domination
$\rho\otimes\rho\le36\mu\otimes\mu$.  A union bound for the two possible
color changes proves Equation~\eqref{eq:variance-transfer}.

Combining the claim with Equation~\eqref{eq:putative-small-cut} and the
tether lower bound gives
\[
 T\le36\kappa\nu+\kappa L_{\rm aux}
 \le36\kappa\nu+T/2,
 \qquad\text{hence}\qquad T\le72\kappa\nu.
\]
Moreover $\nu>0$: otherwise these inequalities give $T=0$, then
$L_{\rm aux}=0$, and Equation~\eqref{eq:variance-transfer} gives
$z(1-z)=0$, contrary to the cut being proper and nonempty.
By Equations~\eqref{eq:raw-capacities}--\eqref{eq:capacity-upper},
\[
 \cD(h)\le T\le72\kappa\nu<T_0\nu.
\]
This contradicts Theorem~\ref{thm:score-gap}.  Every nonempty proper cut in
a NO instance therefore has ratio strictly greater than $C'a$.

\subsection{Complexity and conclusion}

Every enumeration above is polynomial in the original formula size; its
exponent may depend on $C$.  The rational compilation has polynomial bit
complexity, the replication creates polynomially many vertices, and all
capacities and the threshold have polynomial binary encodings.  The
construction does not inspect a satisfying assignment or a candidate cut.
It is therefore a polynomial-time reduction with
\[
 \text{YES}:\ \Phi(G)\le a,
 \qquad
 \text{NO}:\ \Phi(G)>C'a.
\]
A factor-$C$ approximation outputs a cut of ratio at most $Ca<C'a$ in the
YES case, whereas every cut has ratio greater than $C'a$ in the NO case.
Testing the output ratio against the rational number $C'a$ decides the
original satisfiability instance.  This proves Theorem~\ref{thm:main}.
